\DocumentMetadata{
  pdfversion=2.0,
  pdfstandard=ua-2,
  lang=en-US,
  tagging=on,
  tagging-setup={math/setup=mathml-SE,tabsorder=structure}
}
\documentclass[11pt,letterpaper]{article}
\usepackage[margin=0.68in,includefoot]{geometry}
\usepackage{newcomputermodern}
\usepackage{amsmath,amscd,mathtools}
\providecommand{\square}{\mdlgwhtsquare}
\usepackage[hidelinks]{hyperref}
\hypersetup{
  pdftitle={Analysis First-Year Exam, January 2020, Part 2},
  pdfauthor={University of Florida Department of Mathematics},
  pdfsubject={Graduate mathematics examination},
  pdfdisplaydoctitle=true,
  bookmarksopen=true
}
\setcounter{secnumdepth}{0}
\setlength{\parindent}{0pt}
\setlength{\parskip}{0.28em}
\setlength{\emergencystretch}{3em}
\setlength{\leftmargini}{2.15em}
\setlength{\leftmarginii}{2.65em}
\setlength{\leftmarginiii}{3em}
\renewcommand{\labelenumi}{\textbf{\theenumi.}}
\pagestyle{empty}
\raggedbottom
\allowdisplaybreaks
\newenvironment{examproblems}[1]{%
  \begin{enumerate}%
  \setcounter{enumi}{#1}%
  \setlength{\topsep}{0.35em}%
  \setlength{\partopsep}{0pt}%
  \setlength{\itemsep}{0.48em}%
  \setlength{\parsep}{0.18em}%
}{\end{enumerate}}
\newenvironment{examsubparts}{%
  \begin{enumerate}%
  \setlength{\topsep}{0.18em}%
  \setlength{\partopsep}{0pt}%
  \setlength{\itemsep}{0.16em}%
  \setlength{\parsep}{0.1em}%
}{\end{enumerate}}
\newenvironment{examinstructions}{%
  \par\begingroup\itshape
}{%
  \par\endgroup\vspace{0.18em}
}
\makeatletter
\renewcommand\section{\@startsection{section}{1}{0pt}%
  {-1.25ex plus -.25ex minus -.1ex}%
  {0.55ex plus .1ex}%
  {\normalfont\large\bfseries}}
\renewcommand{\@maketitle}{%
  \newpage\null\vskip -1.2em%
  {\footnotesize\bfseries
    \href{https://www.ufl.edu/}{University of Florida}, \href{https://math.ufl.edu/}{Department of Mathematics}\par}%
  \vskip 0.18em%
  \tagstructbegin{tag=Title}%
    {\LARGE\bfseries\href{https://gma.math.ufl.edu/exams/analysis-first-year/}{Analysis First-Year Exam}, January 2020, Part 2\par}%
  \tagstructend%
  \vskip 0.35em%
  \tagmcbegin{artifact}%
    \hrule height 1pt%
  \tagmcend%
  \vskip 0.55em%
}
\makeatother
\title{Analysis First-Year Exam, January 2020, Part 2}
\author{}
\date{}
\begin{document}
\maketitle
\thispagestyle{empty}
\begin{examinstructions}
Do exactly 4 problems. Work must be presented in a neat and logical fashion in order to receive credit. Do not leave any gaps. When a theorem is used in a proof it must be precisely stated.
\end{examinstructions}
\begin{examproblems}{0}
\item
Let~\(f\) and \(f_n\) (\(n>0\)) be real-valued functions on \((0,1)\) with \(f_n\to f\). Decide (with proof) the truth/falsity of each of the following statements when the convergence of \(f_n\) to~\(f\) is (a) pointwise (b) uniform:
\begin{examsubparts}
\item[\textnormal{(i)}]
if each \(f_n\) is increasing on \((0,1)\) then so is~\(f\);
\item[\textnormal{(ii)}]
if each \(f_n\) is bounded then so is~\(f\).
\end{examsubparts}
\item
Let \(f:[0,1]\times[0,1]\to\mathbb{R}\) be a continuous function and let \(\epsilon>0\). Prove that there exist finitely many continuous functions \(g_1,\ldots,g_n:[0,1]\to\mathbb{R}\) and \(h_1,\ldots,h_n:[0,1]\to\mathbb{R}\) such that 
\[
\left|f(x,y)-\sum_{j=1}^{n}g_j(x)h_j(y)\right|<\epsilon \quad\text{for all }(x,y)\in[0,1]\times[0,1].
\]
\item
Let \(f_n\) be a sequence of nonnegative measurable functions and suppose there is an \(L^1\) function~\(g\) such that \(f_n\leq g\) for all~\(n\). Prove that 
\[
\limsup_{n\to\infty}\int f_n\leq\int\limsup_{n\to\infty}f_n.
\]
 Give an example to show the conclusion can fail if the hypothesis \(f_n\leq g\) is removed.
\item
Let \((X,\mathcal{M})\) be a measurable space and \(f_n:X\to\mathbb{R}\), \(n=1,2,3,\ldots\) a sequence of measurable functions. Prove that each of the following subsets of~\(X\) is measurable:
\begin{examsubparts}
\item[\textnormal{a)}]
\(\{x\mid f_n(x)\to+\infty\}\);
\item[\textnormal{b)}]
\(\{x\mid f_{n+1}(x)>f_n(x)\text{ for infinitely many }n\}\);
\item[\textnormal{c)}]
\(\{x\mid f_n(x)\text{ is rational for all }n\}\).
\end{examsubparts}
\item
Let \(E\subset[0,1]\) be a Lebesgue measurable set with \(m(E)>0\). Prove that there exists a point \(0<c<1\) such that 
\[
m(E\cap[0,c])=\frac{1}{2}m(E).
\]
\end{examproblems}
\end{document}
